\chapter{The categorical branch and the conditional MLC roots}

\section{What category theory changes}

The categorical development does not replace the analytic problem by a
formal trick.  It performs two more limited tasks:
\begin{enumerate}
\item it packages subsets of the parameter plane and their inclusions as
      objects and morphisms in a topological over-category;
\item it isolates the exact parameter--dynamical existence statement whose
      connected-image consequence is formal.
\end{enumerate}
The first task is elementary category theory in \texttt{TopCat}.  The second
task is still complex dynamics and is not supplied by the category
interfaces.

\section{The over-category of parameter approximations}

\begin{definition}[Approximation over the parameter plane]
\label{def:categorical-approximation}
Let
\[
 \mathsf{Plane}:=\operatorname{TopCat.of}\C.
\]
An approximation is an object of the over-category
\[
 \mathsf{Approx}:=\mathsf{TopCat}/\mathsf{Plane}.
\]
For $S\subseteq\C$, the canonical approximation is
\[
 \mathsf{ofSet}(S):=
 \bigl(\operatorname{TopCat.of}S\longrightarrow\mathsf{Plane}\bigr).
\]
\end{definition}

The implementation is:
\begin{itemize}
\item \lean{Mlc.CategoricalTopologicalApproximation.Approximation};
\item \lean{Mlc.CategoricalTopologicalApproximation.ofSet};
\item \lean{Mlc.CategoricalTopologicalApproximation.setInclusion}.
\end{itemize}
The underlying parameter subset of an arbitrary approximation is its range:
\[
 \operatorname{image}(A):=
 \operatorname{range}(A\to\mathsf{Plane}).
\]

\begin{theorem}[Images and categorical intersections]
\label{thm:categorical-image-intersection}
For approximations $A,B$,
\[
 \operatorname{image}(A\times_{\mathsf{Plane}}B)
 =
 \operatorname{image}(A)\cap\operatorname{image}(B).
\]
For subsets $S,T\subseteq\C$,
\[
 \mathsf{ofSet}(S)\longrightarrow\mathsf{ofSet}(T)
\quad\Longleftrightarrow\quad S\subseteq T.
\]
\end{theorem}

These are theorems
\lean{Categorical.image_intersection} and
\lean{Categorical.factorsThrough_ofSet_iff}.  They are important because
they show that the categorical pullback is not an abstract replacement for
set intersection: its image is exactly the ordinary set intersection.

\section{The finite-orbit tower}

\begin{definition}[Outer and inner orbit approximations]
\label{def:categorical-orbit-tower}
For $N\in\N$ define
\[
 I_N:=\left\{c:\forall n,\ |p_n(c)|\le N\right\},
\qquad
 O_N:=\left\{c:|c|\le2,\ \forall n\le N,\ |p_n(c)|\le2\right\}.
\]
The inner sets are increasing and the outer sets are decreasing.
\end{definition}

The exact Lean declarations are
\lean{Categorical.Mandelbrot.innerOrbitSet} and
\lean{Categorical.Mandelbrot.outerOrbitSet}.  The following statements are
proved:
\begin{align*}
 I_N&\subseteq\M,
 &\bigcup_N I_N&=\M,\\
 \M&\subseteq O_N,
 &\bigcap_N O_N&=\M.
\end{align*}
The corresponding theorem names are
\begin{center}
\lean{Categorical.Mandelbrot.iUnion_innerOrbitSet_eq_set},
\qquad
\lean{Categorical.Mandelbrot.iInter_outerOrbitSet_eq_set}.
\end{center}

\begin{remark}
The outer and inner limits agree as sets for different reasons.  This
agreement does not say that the finite stages are connected, and it does not
imply local connectivity of the limit.  The missing information is
component control, represented in the repository by the buffer and
finite-component interfaces.
\end{remark}

\section{The categorical root}

\begin{definition}[Root input]
\label{def:categorical-root-input}
The current root input is the proposition
\[
\begin{gathered}
\mathsf{BUF}(\M,O):\quad
\forall c\in\M\ \forall\varepsilon>0\
\exists\,0<\delta<r<\varepsilon,\\
\forall N\ \exists L\ge N,\qquad
O_L\cap\overline B(c,\delta)
\subseteq
\operatorname{Comp}_c(O_N\cap\overline B(c,r)).
\end{gathered}
\]
\end{definition}

Its Lean packaging is
\lean{MLC.RootInput}; the field is
\lean{ParameterComponent.MandelbrotUniformOuterBuffer}.
The topological theorem
\lean{ParameterComponent.locallyConnectedSpace_of_uniformOuterBuffer}
proves
\[
 \mathsf{BUF}(\M,O)\Longrightarrow
 \operatorname{LocallyConnectedSpace}(\M).
\]

\begin{theorem}[Categorical compatibility root]
\label{thm:categorical-compatibility-root}
Let
\[
 \mathsf{MLC}_{\mathrm{cat}}
 :=
 \operatorname{LocallyConnectedSpace}
   \bigl(\operatorname{TopCat.of}\M\bigr).
\]
Then
\[
 \mathsf{MLC}_{\mathrm{cat}}
 \Longleftrightarrow
 \operatorname{LocallyConnectedSpace}(\M),
\]
and $\mathsf{BUF}(\M,O)$ implies both sides.
\end{theorem}

The exact declarations are:
\begin{itemize}
\item \lean{Categorical.MLCConjecture};
\item \lean{Categorical.mlcConjecture_iff_set};
\item \lean{Categorical.categorical_mlc_conjecture};
\item \lean{MLC.mlc_conjecture}.
\end{itemize}
The literature-facing wrappers are:
\begin{itemize}
\item \lean{Literature.categoricalMLC_iff_setMLC};
\item \lean{Literature.categoricalMLC_of_rootInput};
\item \lean{Literature.setMLC_of_rootInput}.
\end{itemize}

\section{Finite covers and uniform geometric towers}

The finite-component branch replaces the pointwise buffer by finite
certificate data.  At scale $k$, a certificate supplies finitely many
rational parameter boxes and a stage $T$ covering $O_T$.  Each local box
comes with radii and a component-control implication.  The theorem
\[
 \mathsf{FC}\Longrightarrow\mathsf{BUF}(\M,O)
\Longrightarrow\mathsf{MLC}(\M)
\]
is implemented by
\lean{FiniteComponent.uniformOuterBuffer_of_finiteComponentCriterion}
and
\lean{FiniteComponent.rootInput_of_finiteComponentCriterion}.

The independent uniform-geometric branch assumes a tower of maps
\[
 R_n:D\longrightarrow D,\qquad
 R_n|_{O_{N_n}}=\id,\qquad
 \|R_{n+1}-R_n\|_\infty\le C2^{-n}.
\]
The limit is a retraction $D\to\M$.  The proved implication is
\[
 \mathsf{UG}\Longrightarrow
 \text{continuous retraction }D\to\M
\Longrightarrow\mathsf{MLC}.
\]
The root theorem is
\lean{UniformGeometry.mandelbrot_locallyConnected_of_squareTower}.
Its tower existence is not asserted.

\begin{remark}
The categorical statement
\lean{UniformGeometry.categorical_mlc_of_squareTower}
is a packaging of the same set-level conclusion in \texttt{TopCat}; it is
not evidence that the tower exists.  Likewise,
\lean{FlowInterfaces.CoherentParametrization} and
\lean{FlowInterfaces.TerminalExtension} are explicit construction
obligations, not project axioms.
\end{remark}

\section{The categorical Green-sublevel branch}

For $c\in\M$ and $n\in\N$ define the frozen translate
\[
 U_n(c):=\{c':G_c(c'-c)<2^{-n}\},
\qquad
 V_n(c):=U_n(c)\cap\M.
\]
The un-intersected set $U_n(c)$ is connected.  The current code proves
\[
 \bigcap_n U_n(c)=(c+K_c),
\]
not $\{c\}$.  This is the theorem
\lean{Quadratic.iInter_green_sublevel_translate_eq_translate_filledJulia}.

The desired straddling statement would be
\[
 U_n(c)\not\subseteq\M
\Longrightarrow V_n(c)\text{ is connected}.
\]
Its categorical form is
\lean{GreenSublevelIntersectionCategoricalData}.
The equivalence with the set-level formulation is proved by
\lean{greenSublevelIntersectionCategoricalData_iff}.

\begin{theorem}[Frozen full-sublevel obstruction]
\label{thm:frozen-sublevel-obstruction}
The universal frozen Green-sublevel intersection statement is false.  There
are $c\in\M$ and $n$ for which $V_n(c)$ is disconnected.
\end{theorem}

The machine-checked witness is $c=i$, $n=16$.  The theorem is
\lean{not_greenSublevelIntersectionCategoricalData}; its set-level source is
\lean{GreenCounterexample.not_greenSublevelIntersectionSetData}.
Consequently, the full Green-sublevel categorical datum cannot be cited as
the classical Yoccoz parapuzzle theorem.

\section{Carving, Efimov-style interfaces, and the literature}

\begin{definition}[Categorical carving datum]
\label{def:categorical-carving}
A carving datum at $(c,n)$ is a surjective continuous morphism
\[
 \operatorname{TopCat.of}U_n(c)\longrightarrow
 \operatorname{TopCat.of}V_n(c).
\]
Since the source is connected, its image is connected.
\end{definition}

The exact native theorem is
\lean{Efimov.isConnected_of_topCatSurjectiveMorphism}.
The stronger analytic input
\lean{Efimov.SpaceHolomorphicCarvingData} supplies a differentiable map
with exact image $V_n(c)$ and is converted to
\lean{Efimov.DouadyHubbardYoccozCategoricalCarvingData}.
The existence proposition is
\lean{Efimov.DouadyHubbardYoccozCategoricalTheorem}.

The Efimov-inspired part is intentionally limited by available foundations:
\begin{itemize}
\item \lean{Efimov.DualizablePacmanTower} records rigid monoidal levels and
      adjoint transitions;
\item \lean{Efimov.StrongMittagLefflerData} records stabilization clauses;
\item \lean{Efimov.PacmanKTheory} records a graded additive shadow, not
      algebraic $K$-theory;
\item \lean{Efimov.PacmanRealization} records a realization into the
      topological approximation category.
\end{itemize}
Mathlib does not provide the stable infinity-categorical and continuous
$K$-theoretic machinery used in Efimov's papers.  Therefore the category
interfaces organize a future realization theorem but do not prove the
missing parameter carving.

\begin{conditionalstep}
The correct relation to the literature is:
\[
\begin{gathered}
\text{Douady--Hubbard/Yoccoz phase--parameter theorem}\\
\Downarrow\\
\text{carving or connected shrinking parapuzzle data}\\
\Downarrow\\
\text{topological connected-image theorem in Lean}\\
\Downarrow\\
\text{pointwise MLC and, after uniformization, the global target}.
\end{gathered}
\]
The categorical layer proves the third line.  It does not supply the first
line, and the counterexample proves that the frozen full-sublevel version of
the first line is false.
\end{conditionalstep}

\section{Categorical status ledger}

\begin{center}\small
\begin{tabular}{@{}P{0.30\textwidth}@{\hspace{4pt}}P{0.48\textwidth}@{\hspace{4pt}}P{0.14\textwidth}@{}}
\textbf{Statement} & \textbf{Lean content} & \textbf{Status}\\
\hline
Set/categorical target equivalence &
\lean{Categorical.mlcConjecture_iff_set} &
proved\\
Root input to categorical MLC &
\lean{Categorical.categorical_mlc_conjecture} &
proved conditionally\\
Finite-component bridge &
\lean{FiniteComponent.rootInput_of_finiteComponentCriterion} &
proved conditionally\\
Uniform tower to MLC &
\lean{UniformGeometry.mandelbrot_locallyConnected_of_squareTower} &
proved conditionally\\
Frozen Green intersection &
\lean{not_greenSublevelIntersectionCategoricalData} &
refuted\\
Space-holomorphic carving &
\lean{Efimov.SpaceHolomorphicCarvingData} &
open input\\
\end{tabular}
\end{center}
